\scanpage{113}
\mspage{A 2}

de Mai \inserted{\handwritten{et Juin}} (l'an dernier) qui suit la ``Promenade''. Cette Lettre était destinée à mes ex-élèves et à mes ``amis d'antan'' dans le monde mathématique. Mais elle non plus n'a rien de technique. Elle peut être lue sans problème par tout lecteur qui serait intéressé à apprendre, par un récit ``sur le vif'', les tenants et aboutissants qui m'ontfinalement amené à écrire Récoltes et Semailles. Plus encore que la Promenade, ça te donnera aussi un avant-goût d'une certaine ambiance, dans le ``grand monde'' mathématique. Et aussi (tout comme la Promenade), de mon style d'expression, un peu spécial paraît-il. Et de l'esprit aussi qui s'exprime par ce style - un esprit qui lui non plus n'est pas apprécié par tout le monde.

Dans la Promenade et un peu partout dans Récoltes et Semailles, je parle du t\,r\,a\,v\,a\,i\,l\quad m\,a\,t\,h\,é\,m\,a\,t\,i\,q\,u\,e. C'est un travail que je connais bien et de première main. La plupart des choses que j'en dis sont vraies, sûrement, pour tout travail créateur, tout travail de découverte. C'est vrai tout au moins pour le travail dit ``intellectuel'', celui qui se fait surtout ``par la tête'', et en écrivant. Un tel travail est marqué par l'éclosion et par l'épanouissement d'une c\,o\,m\,p\,r\,é\,h\,e\,n\,s\,i\,o\,n des choses que nous sommes en train de sonder. Mais, pour prendre un exemple au bout opposé, la passion d'amour est, elle aussi, pulsion de découverte. Elle nous ouvre à une connaissance dite ``charnelle'', qui elle aussi se renouvelle, s'épanouit, s'approfondit. Ces deux pulsions - celle qui anime le mathématicien au travail, disons, et celle en l'amante ou en l'amant - sont bien plus proches qu'on ne le soupçonne généralement, ou qu'on n'est disposé à se l'admettre. Je souhaite que les pages de Récoltes et Semailles puissent contribuer à te le faire sentir, dans ton travail et dans ta vie de tous les jours.

Au cours de la Promenade, il sera surtout question du travail mathématique lui-même. J'y reste quasiment muet par contre sur le c\,o\,n\,t\,e\,x\,t\,e où ce travail se place, et sur les m\,o\,t\,i\,v\,a\,t\,i\,o\,n\,s qui jouent en dehors du temps de travail proprement dit. Cela risque de donner de ma personne, ou du mathématicien ou du ``scientifique'' en général, une image flatteuse certes, mais déformée. Genre ``grande et noble passion'', sans correctif d'aucune sorte. Dans la ligne, en somme, du grand ``Mythe de la Science'' (avec S majuscule s'il vous plait !). Le mythe héroïque, ``prométhéen'', dans lequel écrivains et savants sont tombés (et continuent à tomber) à qui mieux mieux. Il n'y a guère que les historiens , peut-être, qui y résistent parfois,à ce mythe si séduisant. La vérité, c'est que dans les motivations ``du scientifique'', qui parfois le poussent à investir sans compter dans son travail, l'ambition et la vanité jouent un rôle aussi important et quasiment universel, que dans toute autre profession. Ça prend des formes plus

% archive sheet number: 112
